\ifdefined\gkver\else\def\gkver{student}\fi
\documentclass[\gkver]{gaokaozhenti}
\begin{document}

\chapter{2008年四川延考}
%% paperType: 理综物理部分

\begin{enumerate}
\item
%% number: 14
%% paperName: 2008年四川延考第 14 题 6 分
%% typeId: 1
%% type: 单选题
%% chapter: 机械振动
%% point: 简谐运动的振幅
%% method: 临界法
%% score: 6
%% degree: 650
%% duplicateId: 0
%% body:
光滑的水平面上叠放有质量分别为 $m$ 和 $\frac{m}{2}$ 的两木块，下方木块与一劲度系数为 $k$ 的弹簧相连，弹簧的另一端固定在墙上，如图所示。已知两木块之间的最大静摩擦力为 $f$，为使这两个木块组成的系统像一个整体一样地振动，系统的最大振幅为\xzanswer{C}

\onepicture[w=3.21cm]{figs/14.pdf}

\fourchoices[answer=C]
{$\frac{f}{k}$}
{$\frac{2f}{k}$}
{$\frac{3f}{k}$}
{$\frac{4f}{k}$}

%% answer: C
%% memo:
\memoanswer{
研究质量为 $\frac{m}{2}$ 的木块，随质量为 $m$ 的木块一起做简谐运动，其最大加速度为 $a=\frac{f}{\frac{m}{2}}=\frac{2f}{m}$，即整体运动的最大加速度为 $a=\frac{2f}{m}$，则弹簧的最大拉力为 $\left(\frac{m}{2}+m\right)a=kx$，设最大振幅为 $x=A$，解得 $A=\frac{\frac{3}{2}m\cdot\frac{2f}{m}}{k}=\frac{3f}{k}$，故 C 正确。
}

\item
%% number: 15
%% paperName: 2008年四川延考第 15 题 6 分
%% typeId: 1
%% type: 单选题
%% chapter: 分子动理论
%% point: 分子间的作用力
%% method: 概念辨析
%% score: 6
%% degree: 450
%% duplicateId: 0
%% body:
下列说法正确的是\xzanswer{A}

\fourchoices[answer=A]
{大量分子能聚集在一起形成液体或固体，说明分子之间存在引力}
{被活塞封闭在汽缸中的气体体积增大时压强一定减小}
{被活塞封闭在汽缸中的气体温度升高时压强一定增大}
{气体压强的大小只与温度和气体分子的总数有关}

%% answer: A
%% memo:
\memoanswer{
分子间存在引力和斥力，大量分子能聚集在一起形成液体或固体，说明分子之间存在引力，故 A 正确；由气体状态方程 $\frac{pV}{T}=C$ 知，体积增大，压强不一定减小，还要看温度的变化，故 B 错误；温度升高，压强不一定增大，还要看体积的变化，故 C 错误；气体的压强在微观上与气体分子的平均动能和分子的密集程度有关，而温度决定分子的平均动能，但还与分子的密集程度有关，不是与分子总数有关，故 D 错误。
}

\item
%% number: 16
%% paperName: 2008年四川延考第 16 题 6 分
%% typeId: 1
%% type: 单选题
%% chapter: 交变电流
%% point: 交变电流的描述
%% method: 对称法
%% score: 6
%% degree: 700
%% duplicateId: 0
%% body:
一气体放电管，当其两电极间的电压超过 $500\sqrt{3}\UV$ 时，就放电而发光。在它发光的情况下逐渐降低电压，要降到 $500\sqrt{2}\UV$ 时才熄灭。放电管两电极不分正负。现有一正弦交流电源，输出电压峰值为 $1000\UV$，频率为 $50\UHz$。若用它给上述放电管供电，则在一小时内放电管实际发光的时间为\xzanswer{B}

\fourchoices[answer=B]
{10 分钟}
{25 分钟}
{30 分钟}
{35 分钟}

%% answer: B
%% memo:
\memoanswer{
由题意该交流电的表达式为 $e=E\sin 2\pi ft=1000\sin 100\pi t$，在前半个周期内，当 $t=\frac{T}{6}$ 时开始发光，$t=\frac{3}{8}T$ 时停止发光，发光时间为 $\Delta t=\frac{3}{8}T-\frac{T}{6}=\frac{5}{24}T$，整个周期发光时间为 $2\Delta t=\frac{5}{12}T$，故一个小时内的发光时间为 $t=\frac{3600}{T}\times\frac{5}{12}T=1500\Us=25$ 分钟，故 A、C、D 错误。
}

\item
%% number: 17
%% paperName: 2008年四川延考第 17 题 6 分
%% typeId: 2
%% type: 多选题
%% chapter: 静电场
%% point: 电场强度与电势能
%% method: 动力学分析
%% score: 6
%% degree: 650
%% duplicateId: 0
%% body:
如图，在真空中一条竖直向下的电场线上有两点 a 和 b。一带电质点在 a 处由静止释放后沿电场线向上运动，到达 b 点时速度恰好为零。则下面说法正确的是\xzanswer{AC}

\onepicture[align=right, h=3.10cm]{figs/17.pdf}

\fourchoices[answer=AC]
{a 点的电场强度大于 b 点的电场强度}
{质点在 b 点所受到的合力一定为零}
{带电质点在 a 点的电势能大于在 b 点的电势能}
{a 点的电势高于 b 点的电势}

%% answer: AC
%% memo:
\memoanswer{
由于质点运动过程中初末速度均为零，因此质点所受电场力向上，而且先加速后减速，故在 a 点电场力大于重力，b 点重力大于电场力，在 a、b 之间某点重力等于电场力，合外力为零，且此时速度最大，故 A 正确、B 错误；从 a 运动到 b 过程中，电场力做正功，电势能减小，所以 a 点的电势能大于在 b 点的电势能，故 C 正确；沿电场线方向电势降低，a 点的电势低于 b 点的电势，故 D 错误。
}

\item
%% number: 18
%% paperName: 2008年四川延考第 18 题 6 分
%% typeId: 1
%% type: 单选题
%% chapter: 万有引力与航天
%% point: 人造卫星的运行规律
%% method: 比较法
%% score: 6
%% degree: 650
%% duplicateId: 0
%% body:
如图，地球赤道上的山丘 e、近地资源卫星 p 和同步通信卫星 q 均在赤道平面上绕地心做匀速圆周运动。设 e、p、q 的圆周运动速率分别为 $v_{1}$、$v_{2}$、$v_{3}$，向心加速度分别为 $a_{1}$、$a_{2}$、$a_{3}$，则\xzanswer{D}

\onepicture[w=3.46cm]{figs/18.pdf}

\fourchoices[answer=D]
{$v_{1}>v_{2}>v_{3}$}
{$v_{1}<v_{2}<v_{3}$}
{$a_{1}>a_{2}>a_{3}$}
{$a_{1}<a_{3}<a_{2}$}

%% answer: D
%% memo:
\memoanswer{
对于卫星，由万有引力提供向心力得 $G\frac{Mm}{R^{2}}=m\frac{v^{2}}{R}$，解得 $v=\sqrt{\frac{GM}{R}}$，故卫星的轨道半径 R 越大，线速度 v 越小。由于近地资源卫星 p 的轨道半径小于同步通信卫星 q 的轨道半径，故同步卫星 q 的线速度 $v_{3}$ 小于近地资源卫星 p 的线速度 $v_{2}$，即 $v_{3}<v_{2}$；由于同步通信卫星 q 和赤道上的山丘 e 的角速度相同，即 $\omega_{e}=\omega_{q}$，到地心的距离 $R_{q}>R_{e}$，由 $v=\omega R$ 可得 $v_{1}=R_{e}\omega_{e}$，$v_{3}=R_{q}\omega_{q}$，即 $v_{3}>v_{1}$，得到 $v_{2}>v_{3}>v_{1}$，故 A、B 错误；对于 p 和 q，有 $\frac{GMm}{R^{2}}=ma$，可得 $a=\frac{GM}{R^{2}}$，由于 $R_{p}<R_{q}$，则 $a_{p}>a_{q}$，即 $a_{2}>a_{3}$，由 $a=\omega^{2}R$，由于 $R_{q}>R_{e}$，可得 $a_{q}>a_{e}$，即 $a_{3}>a_{1}$，故 $a_{2}>a_{3}>a_{1}$，故 C 错误、D 正确。
}

\item
%% number: 19
%% paperName: 2008年四川延考第 19 题 6 分
%% typeId: 1
%% type: 单选题
%% chapter: 磁场
%% point: 带电粒子在匀强磁场中的运动
%% method: 几何法
%% score: 6
%% degree: 800
%% duplicateId: 0
%% body:
在 x 轴上方有垂直于纸面的匀强磁场，同一种带电粒子从 O 点射入磁场。当入射方向与 x 轴的夹角 $\alpha=45^{\circ}$ 时，速度为 $v_{1}$、$v_{2}$ 的两个粒子分别从 a、b 两点射出磁场，如图所示。当 $\alpha$ 为 $60^{\circ}$ 时，为了使粒子从 ab 的中点 c 射出磁场，则速度应为\xzanswer{D}

\onepicture[w=4.59cm]{figs/19.pdf}

\fourchoices[answer=D]
{$\frac{1}{2}(v_{1}+v_{2})$}
{$\frac{\sqrt{2}}{2}(v_{1}+v_{2})$}
{$\frac{\sqrt{3}}{3}(v_{1}+v_{2})$}
{$\frac{\sqrt{6}}{6}(v_{1}+v_{2})$}

%% answer: D
%% memo:
\memoanswer{
设 a、b、c 三点的坐标分别为 $x_{1}$、$x_{2}$、$x_{3}$，画出粒子的运动轨迹如图所示。

\onepicture[w=6.42cm]{figs/19_an.pdf}

当 $\alpha=45^{\circ}$ 时，速度为 $v_{1}$ 的粒子从 a 点射出磁场，由图中几何关系可知 $x_{1}=\sqrt{2}R_{1}=\frac{\sqrt{2}mv_{1}}{qB}$；速度为 $v_{2}$ 的粒子从 b 点射出磁场，由几何关系知 $x_{2}=\sqrt{2}R_{2}=\frac{\sqrt{2}mv_{2}}{qB}$。当 $\alpha=60^{\circ}$ 时，速度为 $v_{3}$ 的粒子要从 c 点射出磁场，则由几何关系可知 $x_{3}=\sqrt{3}R_{3}=\frac{\sqrt{3}mv_{3}}{qB}$。因 c 为 ab 的中点，故 $2x_{3}=x_{1}+x_{2}$，即 $v_{3}=\frac{\sqrt{6}}{6}(v_{1}+v_{2})$，故 D 正确。
}

\item
%% number: 20
%% paperName: 2008年四川延考第 20 题 6 分
%% typeId: 2
%% type: 多选题
%% chapter: 原子和原子核
%% point: 氢原子能级与跃迁
%% method: 公式法
%% score: 6
%% degree: 700
%% duplicateId: 0
%% body:
用具有一定动能的电子轰击大量处于基态的氢原子，使这些氢原子被激发到量子数为 $n(n>2)$ 的激发态。此时出现的氢光谱中有 N 条谱线，其中波长的最大值为 $\lambda$。现逐渐提高入射电子的动能，当动能达到某一值时，氢光谱中谱线数增加到 $N'$ 条，其中波长的最大值变为 $\lambda'$。下列各式中可能正确的是\xzanswer{AC}

\fourchoices[answer=AC]
{$N'=N+n$}
{$N'=N+n-1$}
{$\lambda'>\lambda$}
{$\lambda'<\lambda$}

%% answer: AC
%% memo:
\memoanswer{
氢原子处于 n 能级向较低激发态或基态跃迁时，可能产生的光谱线条数的计算公式为 $N=C_{n}^{2}=\frac{n(n-1)}{2}$。设氢原子被激发到量子数为 $n'$ 的激发态时出现的氢光谱中有 $N'$ 条谱线，若 $n'=n+1$，则 $N'=\frac{n(n+1)}{2}=N+n$，故 A 正确、B 错误；氢原子能级越高相邻能级差越小，由 $\Delta E=\frac{hc}{\lambda}$，则 $n'>n$，则 $\Delta E'<\Delta E$，所以 $\lambda'>\lambda$，故 C 正确、D 错误。
}

\item
%% number: 21
%% paperName: 2008年四川延考第 21 题 6 分
%% typeId: 1
%% type: 单选题
%% chapter: 相互作用
%% point: 共点力平衡
%% method: 整体法与隔离法
%% score: 6
%% degree: 800
%% duplicateId: 0
%% body:
两个可视为质点的小球 a 和 b，用质量可忽略的刚性细杆相连，放置在一个光滑的半球面内，如图所示。已知小球 a 和 b 的质量之比为 $\sqrt{3}$，细杆长度是球面半径的 $\sqrt{2}$ 倍。两球处于平衡状态时，细杆与水平面的夹角 $\theta$ 是\xzanswer{D}

\onepicture[w=6.00cm]{figs/21.pdf}

\fourchoices[answer=D]
{$45^{\circ}$}
{$30^{\circ}$}
{$22.5^{\circ}$}
{$15^{\circ}$}

%% answer: D
%% memo:
\memoanswer{
设细杆对两球的弹力大小为 T，小球 a、b 的受力情况如图所示，其中球面对两球的弹力方向指向圆心，即有 $\cos\alpha=\frac{\frac{\sqrt{2}R}{2}}{R}=\frac{\sqrt{2}}{2}$，计算得出 $\alpha=45^{\circ}$，故 $F_{Na}$ 的方向为向上偏右，即 $\beta_{1}=\frac{\pi}{2}-45^{\circ}-\theta=45^{\circ}-\theta$，$F_{Nb}$ 的方向为向上偏左，即 $\beta_{2}=\frac{\pi}{2}-(45^{\circ}-\theta)=45^{\circ}+\theta$。两球都受到重力、细杆的弹力和球面的弹力的作用，过 O 作竖直线交 ab 于 c 点，设球面的半径为 R，由相似三角形可得 $\frac{m_{a}g}{OC}=\frac{F_{Na}}{R}$，$\frac{m_{b}g}{OC}=\frac{F_{Nb}}{R}$，计算得出 $F_{Na}=\sqrt{3}F_{Nb}$；取 a、b 及细杆组成的整体为研究对象，由平衡条件得，水平方向上有 $F_{Na}\sin\beta_{1}=F_{Nb}\sin\beta_{2}$，即 $F_{Na}\sin(45^{\circ}-\theta)=F_{Nb}\sin(45^{\circ}+\theta)$，计算得出 $\theta=15^{\circ}$，故 D 正确。

\onepicture[w=5.96cm]{figs/21_an.pdf}
}

\item
%% number: 22
%% paperName: 2008年四川延考第 22 题 12 分
%% typeId: 4
%% type: 实验题
%% chapter: 恒定电流
%% point: 实验：探究弹簧劲度系数与电压的关系
%% method: 转换法
%% score: 12
%% degree: 700
%% duplicateId: 0
%% body:
\begin{enumerate}
	\item
Ⅰ．（6 分）下图中为“双棱镜干涉”实验装置，其中 S 为单色光源，A 为一个顶角略小于 $180^{\circ}$ 的等腰三角形棱镜，P 为光屏。S 位于棱镜对称轴上，屏与棱镜底边平行。调节光路，可在屏上观察到干涉条纹。这是由于光源 S 发出的光经棱镜作用后，相当于在没有棱镜时，两个分别位于图中 $S_{1}$ 和 $S_{2}$ 位置的相干波源所发出的光的叠加。（$S_{1}$ 和 $S_{2}$ 的连线与棱镜底边平行。）

\onepicture[w=6.84cm]{figs/22.pdf}

已知 $S_{1}$ 和 $S_{2}$ 的位置可由其他实验方法确定，类比“双缝干涉测波长”的实验，可以推测出若要利用“双棱镜干涉”测量光源 S 发出的单色光的波长时，需要测量的物理量是 \tkanswer[2.6cm]{$S_{1}$ 与 $S_{2}$ 间的距离}、\tkanswer[2.6cm]{$S_{1}$（或 $S_{2}$）与光屏间的距离} 和 \tkanswer[2.6cm]{干涉条纹间距}。

	\item
Ⅱ．（12 分）下图中所示装置可以用来测量硬弹簧（即劲度系数较大的弹簧）的劲度系数 $k$，电源的电动势为 $E$，内阻可忽略不计；滑动变阻器全长为 $l$，重力加速度为 $g$，V 为理想电压表。当木板上没有放重物时，滑动变阻器的触头位于图中 a 点，此时电压表示数为零。在木板上放置质量为 $m$ 的重物，滑动变阻器的触头随木板一起下移。由电压表的示数 $U$ 及其他给定条件，可计算出弹簧的劲度系数 $k$。

\onepicture[w=6.14cm]{figs/22a.pdf}

\begin{enumerate}
	\item
	写出 $m$、$U$ 与 $k$ 之间所满足的关系式。

	\item
	已知 $E=1.50\UV$，$l=12.0\Ucm$，$g=9.80\Umsq$。测量结果如下表：

	\begin{center}
	\begin{tabular}{|c|c|c|c|c|c|c|}
	\hline
	$m/\Ukg$ & 1.00 & 1.50 & 3.00 & 4.50 & 6.00 & 7.50 \\ \hline
	$U/\UV$ & 0.108 & 0.154 & 0.290 & 0.446 & 0.608 & 0.740 \\ \hline
	\end{tabular}
	\end{center}

	①在给出的坐标纸上利用表中数据描出 $m-U$ 直线。

	\onepicture[w=9.24cm]{figs/22b.pdf}

	②$m-U$ 直线的斜率为 \tkanswer[2cm]{10.0} $\Ukg/\UV$。

	③弹簧的劲度系数 $k=$ \tkanswer[2.5cm]{$1.23\times10^{3}$} $\UNm$（保留 3 位有效数字）。
\end{enumerate}
\end{enumerate}

%% answer:
\jdanswer{
\begin{enumerate}
	\item
$S_{1}$ 与 $S_{2}$ 间的距离；$S_{1}$（或 $S_{2}$）与光屏间的距离；干涉条纹间距。

	\item
（1）$m=\frac{lk}{Eg}U$；
（2）①如图所示；②$10.0$；③$1.23\times10^{3}$。
\end{enumerate}
}

%% memo:
\memoanswer{
光源 S 发出的光经棱镜作用后，相当于在没有棱镜时，两个分别位于图中 $S_{1}$ 和 $S_{2}$ 位置的相干光源所发出的光的叠加。由双缝干涉条纹的间距公式 $\Delta x=\frac{l}{d}\lambda$ 知，$\lambda=\frac{\Delta x}{l}d$，可知需要测量 $S_{1}$ 与 $S_{2}$ 间的距离、$S_{1}$（或 $S_{2}$）与光屏间的距离、干涉条纹间距。

（1）设放置质量 $m$ 的物体时弹簧压缩量为 $x$，由平衡条件得 $mg=kx$，又 $U=\frac{x}{l}E$，解得 $m=\frac{lk}{Eg}U$。

（2）①用描点法作图，作出的图像如图所示。

\onepicture[w=9.24cm]{figs/22b_an.pdf}

②选取 $(0,0)$、$(0.90,9)$ 两点，则斜率为 $\frac{\Delta m}{\Delta U}=\frac{9-0}{0.90-0}\Ukg/\UV=10.0\Ukg/\UV$。

③因 $m=\frac{lk}{Eg}U$，故 $\frac{lk}{Eg}=\frac{\Delta m}{\Delta U}$，$k=\frac{\Delta m}{\Delta U}\cdot\frac{Eg}{l}=10.0\times\frac{1.50\times9.80}{0.120}\UNm\approx1.23\times10^{3}\UNm$。
}

\item
%% number: 23
%% paperName: 2008年四川延考第 23 题 14 分
%% typeId: 6
%% type: 计算题
%% chapter: 磁场
%% point: 安培力与平衡
%% method: 平衡条件
%% score: 14
%% degree: 650
%% duplicateId: 0
%% body:
如图为一电流表的原理示意图。质量为 $m$ 的均质细金属棒 MN 的中点处通过一绝缘挂钩与一竖直悬挂的弹簧相连，弹簧劲度系数为 $k$。在矩形区域 abcd 内有匀强磁场，磁感应强度大小为 $B$，方向垂直纸面向外。与 MN 的右端 N 连接的一绝缘轻指针可指示标尺上的读数，MN 的长度大于 $\overline{ab}$。当 MN 中没有电流通过且处于平衡状态时，MN 与矩形区域的 cd 边重合；当 MN 中有电流通过时，指针示数可表示电流强度。

\onepicture[align=right, w=5.82cm]{figs/23.pdf}

\begin{enumerate}
	\item
	当电流表示数为零时，弹簧伸长多少？（重力加速度为 g）

	\item
	若要电流表正常工作，MN 的哪一端应与电源正极相接？

	\item
	若 $k=2.0\UNm$，$\overline{ab}=0.20\Um$，$\overline{bc}=0.050\Um$，$B=0.20\UT$，此电流表的量程是多少？（不计通电时电流产生的磁场的作用）

	\item
	若将量程扩大 2 倍，磁感应强度应变为多大？
\end{enumerate}

%% answer:
\jdanswer{
\begin{enumerate}
	\item $\frac{mg}{k}$；
	\item M 端接正极；
	\item $2.5\UA$；
	\item $0.10\UT$。
\end{enumerate}
}

%% memo:
\memoanswer{
（1）设弹簧的伸长量为 $\Delta x$，由平衡条件得 $mg=k\Delta x$ ①，由①式得 $\Delta x=\frac{mg}{k}$ ②。

（2）为使电流表正常工作，作用于通有电流的金属棒 MN 的安培力必须向下，因此 M 端应接正极。

（3）设满量程时通过 MN 的电流强度为 $I_{m}$，由平衡条件得 $BI_{m}\overline{ab}+mg=k(\overline{bc}+\Delta x)$ ③，联立①②③并代入数据得 $I_{m}=2.5\UA$ ④。

（4）设量程扩大后，磁感应强度变为 $B'$，由平衡条件得 $2B'I_{m}\overline{ab}+mg=k(\overline{bc}+\Delta x)$ ⑤，由①⑤得 $B'=\frac{k\overline{bc}}{2I_{m}\overline{ab}}$ ⑥，代入数据得 $B'=0.10\UT$ ⑦。
}

\item
%% number: 24
%% paperName: 2008年四川延考第 24 题 18 分
%% typeId: 6
%% type: 计算题
%% chapter: 运动和力
%% point: 整体法与隔离法
%% method: 整体法与隔离法
%% score: 18
%% degree: 700
%% duplicateId: 0
%% body:
水平面上有一带圆弧形凸起的长方形木块 A，木块 A 上的物体 B 用绕过凸起的轻绳与物体 C 相连，B 与凸起之间的绳是水平的。用一水平向左的拉力 F 作用在物体 B 上，恰使物体 A、B、C 保持相对静止，如图。已知物体 A、B、C 的质量均为 m，重力加速度为 g，不计所有的摩擦，则拉力 F 应为多大？

\onepicture[align=right, w=4.90cm]{figs/24.pdf}

%% answer:
\jdanswer{$\sqrt{3}mg$}

%% memo:
\memoanswer{
设绳中张力为 T，A、B、C 共同的加速度为 a，与 C 相连部分的绳与竖直线的夹角为 $\alpha$，对 A、B、C 组成的整体，由牛顿第二定律得 $F=3ma$ ①。

对 B，有 $F-T=ma$ ②。

对 C，有 $T\cos\alpha=mg$ ③，$T\sin\alpha=ma$ ④。

联立①②式得 $T=2ma$ ⑤。

联立③④式得 $T^{2}=m^{2}(a^{2}+g^{2})$ ⑥。

联立⑤⑥式得 $a=\frac{\sqrt{3}}{3}g$ ⑦。

利用①⑦式得 $F=\sqrt{3}mg$ ⑧。

\onepicture[w=5.08cm]{figs/24_an.pdf}
}

\item
%% number: 25
%% paperName: 2008年四川延考第 25 题 22 分
%% typeId: 6
%% type: 计算题
%% chapter: 动量与碰撞
%% point: 多过程运动与碰撞
%% method: 分段分析
%% score: 22
%% degree: 850
%% duplicateId: 0
%% body:
如图，一质量为 $m=1\Ukg$ 的木板静止在光滑水平地面上。开始时，木板右端与墙相距 $L=0.08\Um$；质量为 $m=1\Ukg$ 的小物块以初速度 $v_{0}=2\Ums$ 滑上木板左端。木板长度可保证物块在运动过程中不与墙接触。物块与木板之间的动摩擦因数为 $\mu=0.1$，木板与墙的碰撞是完全弹性的。取 $g=10\Umsq$，求：

\onepicture[align=right, w=5.86cm]{figs/25.pdf}

\begin{enumerate}
	\item
	从物块滑上木板到两者达到共同速度时，木板与墙碰撞的次数及所用的时间；

	\item
	达到共同速度时木板右端与墙之间的距离。
\end{enumerate}

%% answer:
\jdanswer{
\begin{enumerate}
	\item 2 次，$1.8\Us$；
	\item $0.06\Um$。
\end{enumerate}
}

%% memo:
\memoanswer{
（1）物块滑上木板后，在摩擦力作用下，木板从静止开始做匀加速运动。设木板加速度为 a，经历时间 T 后与墙第一次碰撞，碰撞时的速度为 $v_{1}$，由牛顿第二定律和运动学公式得 $\mu mg=ma$ ①，$L=\frac{1}{2}aT^{2}$ ②，$v_{1}=aT$ ③，联立①②③解得 $T=0.4\Us$，$v_{1}=0.4\Ums$ ④。

在物块与木板两者达到共同速度前，在每两次碰撞之间，木板受到物块对它的摩擦力作用而做加速度恒定的运动，因而木板与墙相碰后将返回至初态，所用时间为 T。

设在物块与木板两者达到共同速度 v 前木板共经历 n 次碰撞，则有 $v=v_{0}-(2nT+\Delta t)a=a\Delta t$ ⑤，式中 $\Delta t$ 是碰撞 n 次后木板从起始位置至达到共同速度时所需的时间。⑤式可改写为 $2v=v_{0}-2nT$ ⑥。由于木板的速率只能位于 0 到 $v_{1}$ 之间，故有 $0\leqslant v_{0}-2nT\leqslant2v_{1}$ ⑦，求解上式得 $1.5\leqslant n\leqslant2.5$ ⑧。由于 n 是整数，故 $n=2$。再由①⑤⑧得 $\Delta t=0.2\Us$ ⑨，$v=0.2\Ums$ ⑩。从开始到物块与木板两者达到共同速度所用的时间为 $t=4T+\Delta t=1.8\Us$（11）。

（2）物块与木板达到共同速度时，木板与墙之间的距离为 $s=L-\frac{1}{2}a\Delta t^{2}$（12），联立①②（12）式并代入数据得 $s=0.06\Um$（13）。
}

\end{enumerate}

\end{document}
